\section{方法与决策规则}
\label{sec:method}

我们研究一个冻结模型 $M$ 与固定动作集 $\mathcal{A}$。输入 $x$ 由任务适配器
渲染为状态字符串与带类型问题；模型返回经过模式校验的决策及其概率。策略只
读取第一次调用，并决定是否发出第二次调用。模型权重、问题模板与候选标签集
在搜索过程中从不修改。这样就把推理期控制与在路由目标上微调骨干区分开来。

\begin{figure}[t]
\centering
\begin{tikzpicture}[
  font=\small,
  >=Stealth,
  node distance=5mm,
  box/.style={draw, rounded corners=2pt, align=center, inner sep=3pt,
              minimum height=9mm, fill=white},
  arr/.style={->, thick},
  dashedarr/.style={->, thick, dashed}
]
\node[box] (x) {输入 $x$};
\node[box, right=of x] (m) {冻结模型 $M$};
\node[box, right=of m] (p) {\noul{} 调用\\$p(x)$};
\node[box, right=of p] (g) {带门\\$\ell \le p \le h$};
\node[box, right=of g, yshift=12mm] (c) {\choice{} 调用\\第二视角};
\node[box, right=of g, yshift=-12mm] (n) {\noul{} 标签\\不再调用};
\draw[arr] (x) -- (m);
\draw[arr] (m) -- (p);
\draw[arr] (p) -- (g);
\draw[arr] (g) -- node[above, font=\scriptsize]{$u=1$} (c);
\draw[arr] (g) -- node[below left, font=\scriptsize]{否则} (n);
\node[box, below=22mm of p, fill=gray!8] (fit) {$J(g)=Q(g)-\lambda\,C(g)$};
\node[box, left=of fit, fill=gray!8] (cache) {不可变输出缓存};
\node[box, right=of fit, fill=gray!8] (search) {GA / ES / NSGA-II};
\draw[dashedarr] (cache) -- (fit);
\draw[dashedarr] (fit) -- (search);
\end{tikzpicture}
\caption{路由设定。冻结模型回答第一个 \noul{} 问题；一个两参数门决定是否
向同一模型提出第二个 \choice{} 问题。搜索完全在缓存输出上进行：候选基因组
通过回放缓存评分，优化器只看到标量适应度 $J$。虚线路径离线，实线路径是
部署策略。}
\label{fig:pipeline}
\end{figure}

对二元任务 CodeXGLUE 与 FinDVer，第一个动作是 \noul{}：模型回答一个是/否
问题并返回 $p(x)=P(\text{真}\mid x)$。第二个动作是 \choice{}：模型在两个
命名标签之间选择。在 CodeXGLUE 中，问题是函数是否含有漏洞，以及应标为
\texttt{secure} 还是 \texttt{vulnerable}；在 FinDVer 中，问题是主张是否
被支持，以及应标为 \texttt{supported} 还是 \texttt{not supported}。级联
总是先支付 \noul{}，只有当门打开时才支付 \choice{}。策略是基因组
$g_{\mathrm{cas}}=(\ell,h,u)$，其中 $\ell,h\in[0,1]$，$u\in\{0,1\}$：
\begin{equation}
\pi_{g}(x)=
\begin{cases}
\choice(x), & u=1\!\land\!\ell\le p(x)\le h,\\
\noul(x), & \text{否则}.
\end{cases}
\label{eq:cascade}
\end{equation}
门使用正类概率而不是对称置信度，因为 \noul{} 问题有一个指定的真标签。
基因组修复会把 $\ell$ 与 $h$ 裁剪到 $[0,1]$，在 $\ell>h$ 时交换二者，并把
$u$ 映射到 $\{0,1\}$。两个固定策略是特例：$u=0$ 是仅 \noul{}；
$u=1,\ell=0,h=1$ 是固定仅切换，它仍然支付第一次 \noul{} 调用，因此每行
成本为两次调用。

ProcessBench 的动作空间不同。模型返回候选步骤的排序列表，策略决定是用
top-1 步骤回答还是弃权。基因组 $g_{\mathrm{abs}}=(t,v,m)$，其中
$t,m\in[0,1]$，$v\in\{0,1\}$，使用 top-1 概率 $s_1(x)$ 与 top-1/top-2
间隔 $\Delta(x)$：
\begin{equation}
\mathrm{abstain}(x)=
\begin{cases}
\text{真}, & s_1(x)<t,\\
\text{真}, & v=1\,\land\,\Delta(x)<m,\\
\text{假}, & \text{否则}.
\end{cases}
\label{eq:abstain}
\end{equation}
当所需置信度字段缺失时，评估器按失败关闭处理：该行被强制回答，并计为错误，
而不是被悄悄丢弃。我们报告覆盖率（被回答行的比例）与选择性准确率（被回答
行上的准确率）。两个基因组从不共享动作集或成本模型，因此我们从不合并它们
的准确率数字。

适应度把任务质量与第二次调用的成本结合起来：
\begin{equation}
J(g)=Q(g)-\lambda\,C(g),
\label{eq:fitness}
\end{equation}
其中 $Q$ 在二元任务上是 macro-F1，在 ProcessBench 上是步骤准确率；$C(g)$
是每行模型调用均值减一；$\lambda=0.05$ 在搜索前固定。Macro-F1 作为主指标，
因为两个二元数据集都类别不平衡；准确率同时报告。第一次 \noul{} 调用计入
每个策略，包括仅 \noul{}，因此成本项隔离的是增量调用，而不是总推理账单。
模型输出只生成一次，并从不可变 JSON 缓存回放，这使候选评估确定且廉价。
我们测量原生动作的端到端延迟，但不为每个搜索出的策略声称延迟。

GA 使用锦标赛选择、均匀交叉、高斯变异与 $u$ 的位翻转。ES 是
$(\mu+\lambda)$ 策略，变异两个端点并固定 $u=1$。NSGA-II 维护非支配前沿，
在二元任务上优化准确率与负调用成本，在 ProcessBench 上优化选择性准确率与
覆盖率；由于其质量目标是准确率而不是 macro-F1，我们单独保留其区间。主预算
是 24 次唯一基因组评估，种子 17、23、41；五种子扩展加入 59 与 73。一个
探针把 CodeXGLUE/Vela 的预算提高到 48。随机搜索使用相同预算与缓存。

策略类的一个性质塑造了校准实验。门只通过区间 $\ell\le p\le h$ 依赖
$p(x)$。任何严格单调变换——例如温度缩放——都把区间映射为另一个区间，并在
端点重参数化意义下诱导同一组被切换行。因此，单调重校准无法改变固定池上
可达的最优带策略。它可以移动行在带之间的位置，但无法创造原始排序本就不
包含的“可修复行 vs 有害行”区分。因此校准与路由分开评估。
